\subsection{Ambient cyclotomic traces for quartic dimension fields}
\label{sec:cyclotomic-stabilizer}
The relative-trace method does not require the dimension field to be
quadratic. One may instead embed it into a cyclotomic field and retain
all coefficients of the resulting normalized traces.

\begin{proposition}[A finite ambient trace domain]
\label{prop:ambient-trace-domain}
Let $L=\QQ(\zeta_F)$ contain all simple dimensions of a pseudo-unitary
fusion category, and let $S$ be the prime support supplied by the
categorical Cauchy theorem. Suppose every prime dividing $F$ lies in
$S$. For height $h$, put
\[
 B=\prod_{p\in S}p^{\max\{v_p(F),1\}+\lfloor\log_p h\rfloor}.
\]
Every joint normalized relative-trace vector
\[
 \left(\frac{\operatorname{Tr}_{L(\theta)/L}(\theta^j)}
 {[L(\theta):L]}\right)_{j=1}^{h}
\]
of a root of unity whose order has prime support in $S$ is either zero,
or occurs among roots of unity of order dividing $B$.
\end{proposition}
\begin{proof}
If $p^a$ divides the order with
$a>\max\{v_p(F),1\}+\lfloor\log_p h\rfloor$, then for each
$1\le j\le h$ the $p$-primary part of $\theta^j$ has order strictly
larger than $p^{\max\{v_p(F),1\}}$. Summing over the last degree-$p$
cyclotomic layer gives zero: its conjugates differ by the complete set
of $p$th roots of unity. This layer fixes $L$ and the other primary
components, so transitivity of trace gives zero in $L$. Otherwise the
order divides $B$. Since $F\mid B$, all remaining traces can be computed
inside $\QQ(\zeta_B)$ by averaging the automorphisms
$\zeta_B\mapsto\zeta_B^a$ with $a\equiv1\pmod F$.
\end{proof}

Use each pair consisting of an admissible high restriction column $v$
and a trace vector as a nonnegative count variable. The Gram equations
remain integral. Apply normalized relative trace to the central
indicator formula of \cite[Theorem~4.1]{NS}; the low induction objects
have twist one. Expanding in the power basis of $L$ gives rational
linear equations for the first three moments. The usual second
indicator choices are exhausted. When $3\nmid F$ and $3\notin S$, the
third indicator lies both in $\QQ(\zeta_3)$ and in a cyclotomic field
of conductor coprime to three, hence is rational. Its cyclic
multiplicities give the necessary description
\[
 \nu_3(Y)=a-b/2,\qquad a,b\ge0,\qquad
 a+b=\dim\operatorname{Hom}(\mathbf1,Y^3).
\]
The count $b$ includes both primitive third-root eigenspaces.
For a four-dimensional cubic space with the anticommuting stabilizer
of Proposition~\ref{prop:stabilizer-rotation}, impose in addition
$\nu_3(Y)\le1$. These conditions are linear after adding nonnegative
slack variables. Allowing rational multiplicities enlarges the
necessary system and is therefore legitimate for exclusion.

